\section{Passage moments and angular transfer}\label{sec:angular}

The actual-path disk estimate of Theorem~\ref{thm:actual-disk} first
gives all exponential moments for raw passages. These moments control
local nest weights and yield a strict cost gap at a flat change of wire
color, at the fixed target parameter $0<q<1$.

The second part isolates two angular decay powers in finite-circle
transfer matrices. This calculation is valid throughout $0<q<4$ and
uses no geometric probability estimate or infinite-volume law. Its
physical interpretation, together with the nest costs, will supply the
current and charge estimates in Section~\ref{sec:physical}.

\subsection{Raw passage moments}

The estimates in this subsection concern the finite switch drawing:
raw strands follow switches and stop at wall ports.  In particular,
neither a strand nor an actual crossing may use an exterior
identification as an edge.  All conditional laws below are obtained by
fixing exterior states in one of the regular buffered geometries of
Section~\ref{sec:disks}; their induced compatible boundary partitions
are allowed to vary.  In every comparison of exterior completions,
the local patch, its fixed buffer, and the continuous designated wire
portions on its walls agree.  Exterior data have the permitted planar
realizations of Definition~\ref{disk:regular-collars}.  For the bounded
collar comparisons used here, a fixed number $J_{\mathrm{bank}}$ bounds
every designated bank-wire class meeting a collar, including a class that
ends inside it.
After those classes are removed, the residual rim caps are separately
noncrossing.  No other identification bypasses the buffer.  Constants may
depend on this fixed inventory of bank classes.

\begin{convention}[Fixed geometry before small mesh]
\label{angular:mesh-convention}
A macroscopic annulus, sector, corridor, or box comes with a fixed
buffer.  All subdivision numbers, aspect ratios, and separations used
in its proof are chosen before the mesh tends to zero.  A bound is
uniform once the mesh resolves these fixed choices.  The remaining
bounded range of radius-to-mesh ratios is covered by the deterministic
bound on the number of tiles in the local patch.  Thus a constant may
depend on a fixed exponential rate or subdivision number, but never
on the subsequently refined mesh.
\end{convention}

\begin{definition}[Raw passage packing]\label{angular:passage-packing}
Fix a finite test patch with its cut-surface drawing.  For an annulus or
sector test, an admissible subarc is contained in that patch and joins its
two prescribed rims.  It is trimmed at its first arrival at the second rim
and the last departure from the first rim before that arrival.  For a diameter test,
fix a positive macroscopic threshold and require instead a contained
raw subarc whose diameter is at least that threshold.  Raw strands stop
at wall ports in both tests.

The passage count of the test is the maximum cardinality of a family of
admissible subarcs that are pairwise disjoint, including their port
incidences on the cut surface.  A finite-family passage count is the sum
of the counts of its nominated tests.  We also allow an inventory bounded
by a fixed multiple of such a sum plus a fixed constant, recording that
domination when it is used.  The count depends only on switches in its
test patches; their larger buffers are used for conditional estimates.
\end{definition}

\begin{proposition}[All exponential passage moments]
\label{prop:passage-moments}
For every fixed buffered regular annulus or sector and every
\(B\ge0\), its passage count \(X\) satisfies
\begin{equation}\label{angular:passage-moment}
 \EE\!\left[e^{BX}\mid\text{exterior states}\right]\le C_B.
\end{equation}
The estimate is uniform in scale, sufficiently fine mesh, and the
compatible exterior states.  It also holds for the fixed macroscopic
diameter tests of Definition~\ref{angular:passage-packing}, for a fixed
finite family of such tests, and for an inventory with the stated fixed
domination by their packing counts.
\end{proposition}

\begin{proof}
\textbf{An ordered crossing test gives one exponential rate.}\quad
We first prove a geometric tail with one positive exponential rate.
For each configuration, select a maximizing packing only to derive the
following deterministic implications.  No conditional law below fixes
that choice of packing.
Cover the relevant buffered annular slice by finitely many small
bulk boxes and boxes flush with a straight bank.  A subarc moving a
fixed fraction of the slice width has a subarc between two separated
tile-coordinate levels in one of these boxes.  This follows by
subdividing each coordinate into a fixed finite grid of levels and
trimming at a last departure and a first arrival.  Slight cropping
keeps the levels inside the available buffer.  Near a corner or a
tip, make the selections in boxes away from the individual apex;
an arc of the prescribed diameter must travel through one of them.
There are only finitely many choices of box, coordinate, and pair
of levels. The same selection applies to a diameter test: an arc of
the prescribed diameter has a coordinate displacement bounded below
by a fixed fraction of that diameter. Thus the argument below proves
the first exponential rate for both crossing and diameter counts.

Suppose one choice receives many disjoint raw subarcs.  They are
ordered separators in the rectangular strip between its two levels.
Small outward prolongations at lateral boundary contacts permit the
usual disk order without adding any random tiles.  On either side
of a raw separator, its incident corner sites form an actual path:
at each switch consecutive sites are either equal or joined by an
actual open bond.  One side has one color and the other side has
the other color.  These paths do not cross a raw separator.
By retaining a fixed fraction of the separators and leaving unused
separators between selections, we obtain disjoint actual crossings
whose colors alternate in strip order.  Trim them to be simple and
to have no end-level contacts except their initial and terminal
vertices.  Lateral contacts are harmless.  Their initial points
retain their order along the first level.

Partition that level into a fixed number of intervals of length
at most \(\eta\) times the level separation.  One interval contains
a consecutive block of at least a fixed fraction of the alternating
crossings, up to a fixed additive loss.  We choose the fixed
\(\eta>0\) sufficiently small below.  Thus a large passage count
forces a long ordered alternation of proper actual crossings
starting in one such interval.

We now test this implication without conditioning on the existence
of the entire alternation.  For the first color, perform complementary
path/cut searches from both lateral sides of the strip, looking for
proper actual crossings from the chosen interval to the other level.
For these searches only, ignore bonds running along the two end
levels and split the incidences at an end-level site into separate
leaves in their fan order.  Give the target leaves auxiliary outward
spokes.  A simple path in this split drawing is exactly a physical
crossing with no additional end contacts.  No switch state or
probabilistic boundary identification has been changed.

Seed each search at its prescribed exterior flank face.  An interior
face is entered only from an already reached face, across a bond that
has been queried and found closed to the tested color.  The search
therefore cannot inspect strictly beyond a witnessing crossing.
Equivalently, after
a successful search, closing every still uninspected tested-color
edge leaves a crossing by the planar path/cut alternative, so one
may select a simple blocker made of measured edges.  An existing
proper crossing shields all edges strictly on its other side.
Consequently, when a long ordered alternation exists, the two
searches find disjoint measured blockers, with all but a bounded
number of the intervening witnesses still between them.

Condition on the two search leaves.  The disk between these
blockers is uninspected in its interior.  For the opposite color,
its longitudinal flanks are free, and its possible inherited
identifications occur only at the two end cuts.  To see this at
tile level, each cutter follows actual open diagonals of the
blocker color.  The opposite-color half-tile slots along it are
dead or dangling slots, not entrances to an exterior wire.
Splitting a lateral contact into its fan incidences does not
introduce a longitudinal connection.  In the conditional
opposite-color graph, all bonds crossing a cutter are closed;
outside connections therefore meet the intercut disk only at
its end cuts.  Unqueried exterior states can be integrated out
after this observation.

The near end of this disk remains within the chosen interval,
whereas the other end is at the original positive level
separation.  Theorem~\ref{thm:actual-disk}, including its allowed
end-cap identifications, gives
\begin{equation}\label{angular:one-round}
 \PP(\text{an opposite-color proper traversal}
       \mid\text{the two search leaves})\le p_0<1
\end{equation}
once \(\eta\) is sufficiently small and the mesh is sufficiently
fine.  The bound is uniform over the irregular measured blockers.

Repeat the two lateral searches inside the current disk, now
testing the opposite color, and then alternate colors.  At each
stage the first and last witnesses of the tested color shield
the remaining middle witnesses; discarding at most two witnesses
at each end preserves the claim.  New blockers cannot use dead
longitudinal slots or the original lateral wires.  The searches
stop when the required disjoint pair is absent.  A long initial
alternation forces linearly many successful rounds.  Except for
the first round, success requires the event in
\eqref{angular:one-round}.  Iterated conditioning on the search
leaves, not conditioning on the desired alternation, therefore
gives
\[
 \PP(X\ge m\mid\text{exterior states})
       \le C e^{-\beta m}
\]
for some fixed \(\beta>0\).  Taking the finite union over the
initial geometric choices proves \eqref{angular:passage-moment}
at one rate \(B_0>0\).

\paragraph{Separated slices give every fixed exponential rate.}
To obtain every fixed rate, divide the radial distance into \(J\)
separated slices of thickness comparable to \(r/J\), with \(J\)
fixed.  A middle rim of each slice has a cover by \(O(J)\) small
boxes whose buffers remain inside that slice.  Use boxes flush
with a wall where necessary, treating the two banks separately.
Their overlap graph has a coloring with \(h\) colors, where \(h\)
is independent of \(J\).  The box shapes range over a fixed family
after rescaling, so their passage counts \(X_\alpha\) have the same
first rate \(B_0\).  For each selected original passage and each slice,
choose a point where it meets the middle rim and a box whose smaller core
contains that point.  Following the passage toward either rim of the slice
until its first exit from the box gives a contained subarc of the box's
fixed diameter threshold;
the box buffer remains in the slice.  The selected original subarcs are
disjoint at all cut-surface incidences, so these contributions are
disjoint as well.  Hence
\[
 JX\le\sum_\alpha X_\alpha .
\]
Within one color class the buffered supports are disjoint.
Conditional iteration of the first-rate estimate gives
\(\EE\exp(B_0\sum_{\alpha\in\mathcal C}X_\alpha)
\le C^{|\mathcal C|}\).
H\"older's inequality over the \(h\) classes now gives a bound at
rate \(B_0J/h\) for \(X\).  Choosing \(J\ge hB/B_0\) proves the
desired rate.  Convention~\ref{angular:mesh-convention} handles
the finite range before the chosen boxes fit in the mesh.
To transfer every rate to a diameter threshold $a>0$, choose a finite
cover of the test patch by smaller neighborhoods inside fixed regular
bulk, wall, or corner charts. Take each chart's outer diameter less
than $a/2$, with its passage-test buffer inside the original available
buffer. Every point belongs to an inner neighborhood separated from
the chart's outer rim by a fixed positive distance. An arc of diameter
at least $a$ therefore has a subarc crossing from such an inner
neighborhood to that outer rim. These are a fixed finite family of
annulus or sector tests. Assign each member of a disjoint diameter
packing to one such crossing; within each test the assigned subarcs
remain disjoint. The diameter count is bounded by the sum of these
crossing counts. Their every-rate bounds and H\"older's inequality
give the assertion. The same inequality combines any fixed finite
family and any inventory dominated by its packing sum.
\end{proof}

\subsection{Nest weights under local tilts}

The passage bounds now control the change in a positive weight when its
nest cutoff moves through one scale.  This comparison must hold under
the weighted law itself, since that law will enter the charge estimates.

Let \(r_j=2^jr_0\), starting at a fixed multiple \(r_0\) of the
mesh.  Let \(N_j\) count raw nests wholly contained in the
corresponding radius-\(r_j\) neighborhood.  A nest is either a
closed raw loop surrounding a marked bulk site or a raw arc
joining the two side rays around a marked boundary gap, corner,
wire change, or slit tip.  These definitions depend on the raw
trace, not on any exterior completion.  The initial count is
bounded deterministically.  There is a fixed integer \(k\) and
shell passage counts \(X_j\) such that
\begin{equation}\label{angular:nest-increments}
 0\le N_j-N_{j-1}\le X_j,
\end{equation}
and the buffered support of \(X_j\) is disjoint from
\(B_{r_{j-k}}\).  Indeed a newly counted nest either travels a
fixed fraction of the radius or winds across the local sector
within that band.  At a slit tip, a route between banks away from
the tip must cross the full-angle sector.  A fixed small-box
cover away from the marked point gives the required counts.
Increase \(k\) once to accommodate all these buffers.

\begin{proposition}[Local tilt bounds]\label{prop:tilts}
Fix \(b>0\), and write
\[
 F_j(b)=\EE b^{N_j},\qquad
 \dd\nu_j=F_j(b)^{-1}b^{N_j}\,\dd\PP.
\]
There is a constant \(C_b\) such that
\begin{equation}\label{angular:tilt-ratios}
 C_b^{-1}\le F_j(b)/F_{j-1}(b)\le C_b .
\end{equation}
For every fixed \(B\ge0\),
\begin{equation}\label{angular:tilt-moments}
 \sup_{i\ge j}\EE_{\nu_j}e^{BX_i}\le C_{b,B}.
\end{equation}
The assertions hold conditionally on compatible exterior states
and allow a fixed additional family of passage counts whose
buffers lie at or beyond the tilt level.  They also hold with
the radial direction reversed, for nests whose raw traces lie
outside a moving inner cutoff and inside a fixed outer cutoff.
\end{proposition}

\begin{proof}
For \(b=1\), this is Proposition~\ref{prop:passage-moments}.
\paragraph{Weights below one.}
Suppose \(0<b<1\), put \(L=|\log b|\), and set
\(D_i=\log(F_{i-1}(b)/F_i(b))\ge0\).
Jensen's inequality and \eqref{angular:nest-increments} imply,
for every fixed \(\varepsilon>0\),
\[
 D_i\le L\EE_{\nu_{i-1}}X_i
 \le\varepsilon\log
       \EE_{\nu_{i-1}}e^{LX_i/\varepsilon}.
\]
Since \(b^{N_{i-1}}\le b^{N_{i-k}}\), and the latter tilt
is supported outside the buffer of \(X_i\), conditioning and
Proposition~\ref{prop:passage-moments} give
\[
 \EE_{\nu_{i-1}}e^{LX_i/\varepsilon}
 \le C_{\varepsilon}\frac{F_{i-k}(b)}{F_{i-1}(b)}.
\]
Thus, changing the fixed constant,
\begin{equation}\label{angular:decreasing-tilt-induction}
 D_i\le C_{\varepsilon}
          +\varepsilon\sum_{i-k<\ell<i}D_\ell.
\end{equation}
The bounded initial count supplies the initial bounds.  Choose
\(\varepsilon(k-1)<1\); induction then bounds every \(D_i\).
For \(i\ge j\), either the buffer of \(X_i\) misses the entire
tilt, or dropping at most \(k\) decreasing increments reduces
to that situation.  The already bounded ratios pay for these
increments and prove \eqref{angular:tilt-moments} at every
fixed rate.

\paragraph{Weights above one.}
Now suppose \(b>1\), and write \(L=\log b\).  Given a desired
rate, choose \(B>kL\) at least that large.  We induct on \(j\)
that
\[
 \EE_{\nu_j}e^{BX_i}\le C_*\qquad(i\ge j).
\]
If \(i\ge j+k\), the tilt and buffer are disjoint, and the
conditional passage estimate applies directly.  Otherwise put
\(j_0=i-k<j\), treating levels below the initial one as no tilt
and paying a fixed factor for the bounded initial count.
Because \(F_j(b)/F_{j_0}(b)\ge1\),
\[
 \EE_{\nu_j}e^{BX_i}
 \le
 \EE_{\nu_{j_0}}\exp\left(
       BX_i+L\sum_{j_0<\ell\le j}X_\ell\right).
\]
There are at most \(k\) summands in the sum.  Give each its
H\"older weight \(L/B\).  Its resulting moment is
\(\EE_{\nu_{j_0}}e^{BX_\ell}\le C_*\), by the induction
hypothesis at the earlier tilt \(j_0\).  The remaining H\"older
weight is at least \(1-kL/B>0\).  Its factor is a moment of
\(X_i\) at a larger fixed rate, under a tilt whose support
misses the buffer of \(X_i\).  Consequently
\[
 \EE_{\nu_j}e^{BX_i}\le C C_*^{kL/B}.
\]
Since \(kL/B<1\), a sufficiently large fixed \(C_*\) closes
the induction.  This proves every desired rate.  The ratio
bound follows from
\[
 1\le F_j(b)/F_{j-1}(b)
    =\EE_{\nu_{j-1}}b^{N_j-N_{j-1}}
    \le\EE_{\nu_{j-1}}e^{LX_j}.
\]
The preceding proofs are conditional and remain valid after
adjoining finitely many passage counts at the indicated levels.

\paragraph{Moving the inner cutoff.}
For the reverse direction, lower the inner cutoff while keeping
the outer one fixed, counting nests whose entire raw traces lie
outside the moving cutoff.  Every newly included nest supplies
a passage near that cutoff by the same radial or sector-crossing
argument used in \eqref{angular:nest-increments}.  A fixed number
of reversed levels separates its buffer from the earlier,
farther-out tilt.  The two inductions therefore apply verbatim
with the radial order reversed.
\end{proof}

\begin{corollary}[Polynomial costs and fixed cutoffs]
\label{angular:polynomial-costs}
Up to radii \(O(n^C)\) in mesh units, every fixed positive
nest weight and each of its fixed real powers has expectation
between \(n^{-A}/C'\) and \(C'n^A\), for constants depending
on the fixed weight and geometry.  Changing a cutoff by a
fixed factor changes its expectation by a bounded factor.
Changing a permitted exterior completion beyond such a fixed
buffer, while keeping the patch, buffer, and deterministic wall
portions unchanged, does the same.  For a fixed weight \(W\), uniformly
for \(t\) in a fixed compact interval,
\[
 \left|\log\EE W^{t+\varepsilon}
              -\log\EE W^t\right|
       \le C|\varepsilon|\log n+O(|\varepsilon|)
\]
when \(\varepsilon\) stays in a sufficiently small fixed
interval.
\end{corollary}

\begin{proof}
There are \(O(\log n)\) dyadic levels, so
\eqref{angular:tilt-ratios} proves the polynomial bounds and
the fixed-factor cutoff assertion.  The bounded reciprocal
marginal-density comparisons in Lemmas~\ref{disk:initial-collars}
and~\ref{disk:thin-collars} give
the exterior assertion.  Finally \(t\mapsto\log\EE W^t\)
is convex.  Its values on a slightly larger compact interval
are \(O(\log n)\), by the bounds just proved; its secant
slopes on the smaller interval are therefore \(O(\log n)\).
\end{proof}

\begin{lemma}[Absorption of macroscopic passage errors]
\label{angular:absorption}
Suppose \(Y_n\ge0\),
\(\EE Y_n\ge C^{-1}n^{-A}\), and
\(\EE Y_n^2\le Cn^A\).  Suppose \(X_n\ge0\) has uniformly
bounded exponential moments at every fixed rate.  For fixed
\(C_0\ge1\) and every \(\varepsilon>0\),
\[
 C_\varepsilon^{-1}n^{-\varepsilon}
 \le
 \frac{\EE[Y_n C_0^{\pm X_n}]}{\EE Y_n}
 \le C_\varepsilon n^\varepsilon.
\]
\end{lemma}

\begin{proof}
Put \(L=\log C_0\) and choose \(\eta>0\) with
\(\eta L<\varepsilon/2\); the case \(L=0\) is immediate.
On \(X_n\le\eta\log n\), either multiplier lies between
\(n^{-\varepsilon/2}\) and \(n^{\varepsilon/2}\).
For every fixed \(t>0\), Cauchy--Schwarz gives
\[
 \EE[Y_ne^{LX_n};X_n>\eta\log n]
 \le (\EE Y_n^2)^{1/2}
      \bigl(n^{-t\eta}\EE e^{(2L+t)X_n}\bigr)^{1/2}
 \le C_t n^{A/2-t\eta/2}.
\]
The same bound, with \(L=0\), controls the omitted
\(\EE[Y_n;X_n>\eta\log n]\).  Choose \(t\) so large that
both errors, after division by \(\EE Y_n\), are smaller
than any prescribed inverse power.  Retaining the truncated
event proves the lower bounds, and adding the tail proves
the upper bounds.  Enlarge constants for bounded \(n\).
\end{proof}

\subsection{The strict cost at a flat wire change}

Consider a straight boundary with one change of wire color,
a marked gap at that change, and a fixed buffered half-box.
The gap is the boundary incidence of its designated physical tile
corner between neighboring wall ports.  Its own color is the primal
or dual color of that corner, independently of the wire designations
on the two adjacent wall intervals.  Distinct incidences at a cut
remain distinct, and no wall cap is used to extend an actual component.
Let \(\mathcal A_n\) be the event that the marked gap has an
actual arm of its own color to radius comparable to \(n\)
in mesh units.  Let \(N_n\) count raw arcs straddling the
gap, with one endpoint on each side ray, whose entire traces
are inside that cutoff.  Define
\begin{equation}\label{angular:flat-cost-definitions}
 P(n)=\log\PP(\mathcal A_n),\qquad
 D_*(n)=\log\EE d^{N_n},\qquad d=\sqrt q.
\end{equation}
Here \(0<d<1\), so no nonnegative sign is asserted for \(D_*\).
All definitions use a fixed compatible positive exterior
completion beyond the buffer.

\begin{lemma}[Flat costs, gap arms, and raw straddles]
\label{angular:arm-straddle}
\label{angular:flat-costs}
There are fixed \(0<c<C<\infty\) such that, with harmless
fixed changes of the cutoff convention,
\begin{equation}\label{angular:arm-straddle-inclusions}
 \mathcal A_{Cn}\subseteq\{N_n=0\}\subseteq\mathcal A_{cn}.
\end{equation}
Arm probabilities at fixed multiples of a radius are
comparable.  Moreover
\[
-A\log n-O(1)\le P(n)\le-a\log n+O(1)
\]
for some \(A,a>0\), and
\[
 -A\log n-O(1)\le D_*(n)\le0.
\]
Changing a cutoff by a fixed factor or changing a compatible
exterior completion beyond a fixed buffer changes either
\(P(n)\) or \(D_*(n)\) by \(O(1)\).
These statements are uniform under the compatible exterior
changes allowed above.
\end{lemma}

\begin{proof}
Work in the cut drawing without using exterior caps or wire
identifications in its connectivity test; the sampling law retains
its prescribed boundary partition.  Within each tile, the corner sites
on a given side of a switch strand are equal or joined by an actual
open diagonal.  Concatenating these local incidences identifies
boundary-touching regions of the drawing with the actual color
components incident to the boundary gaps.  This identification uses
only physical edges.  The disjoint raw boundary-to-boundary arcs
have a tree as their region-adjacency graph.  A closed raw loop has
no wall port and encloses no boundary gap; it therefore does not
split a boundary-touching class.  Apparent bank contacts in an
unsmoothed drawing retain their distinct fan incidences on the cut
surface and do not join such classes.

Let $S$ be the set of tree vertices visited by the outer-rim walk.
It is a connected subtree; neither that walk nor the wall walk is
assumed to be a simple path.  If the region of the marked gap is
absent from $S$, take the first edge separating it from $S$.
An arc with an endpoint on the rim has both adjacent regions
visited by the rim walk, so this separating edge cannot have such
an endpoint.  Both ends of its raw arc are on the wall.  The wall
walk enters and leaves the component containing the gap across
that edge, hence the two ends lie on opposite sides of the gap.
This is a straddling raw arc.  Conversely, a wall-to-wall raw arc
straddling the gap separates it from a sufficiently distant rim,
and an actual path of the gap color cannot cross it.  The region
argument is unchanged by interior closed loops or mixed wire
designations, since no exterior identification has been used.
Applying it in two nested buffered half-boxes, with fixed margins
to absorb the tile diameter, proves
\eqref{angular:arm-straddle-inclusions}.

For a fixed-factor extension use the single-arm version of
Lemma~\ref{lem:relative-arm}.  Its finitely many truncated-search
success events $T_1,\ldots,T_m$ cover the shorter-arm event.  On
each successful leaf, Proposition~\ref{prop:routing} and
Lemma~\ref{lem:confined-join} extend the measured source-connected
frontier through the remaining window with probability at least
$c_0$.  Thus, for the permitted states $F$ fixed away from that
window,
\[
 \PP(\mathcal A_{Cn}\mid F)
 \ge c_0\max_i\PP(T_i\mid F)
 \ge \frac{c_0}{m}\PP(\mathcal A_n\mid F).
\]
This is a relative event estimate, not an assertion under arbitrary
conditioning on the full states of an incoming arm.  Iterating from a
bounded-mesh event of positive finite-energy probability
gives the polynomial lower bound.  In the other direction,
Proposition~\ref{prop:screens} provides opposite-color
actual wall-to-wall screens in separated annuli, each with
uniformly positive conditional probability.  Each screen
blocks the arm.  Conditional iteration over \(O(\log n)\)
annuli gives the strict polynomial upper bound.
The bounded collar comparisons in Lemmas~\ref{disk:initial-collars}
and~\ref{disk:thin-collars} make these estimates uniform under
the specified exterior changes.  None of
these paths traverses a wire identification.
The assertions for \(D_*\) follow from
Corollary~\ref{angular:polynomial-costs}; its upper bound
is immediate from \(d<1\) and \(N_n\ge0\).
\end{proof}

\begin{lemma}[Strict flat-flip weight gap]\label{lem:strict-flip}
For every fixed \(0<b_1<b_2\), there is \(c>0\) such that
\begin{equation}\label{angular:strict-weight-gap}
 \log\EE b_2^{N_n}-\log\EE b_1^{N_n}
       \ge c\log n-O(1).
\end{equation}
In particular, for some \(c'>0\),
\begin{equation}\label{angular:strict-flat-gap}
 D_*(n)\ge P(n)+c'\log n-O(1).
\end{equation}
\end{lemma}

\begin{proof}
Fix \(b=b_1\).  At a dyadic scale \(r_i\) well between
the mesh and the cutoff, choose a fixed large integer
\(k\), larger than all buffer ranges in
Proposition~\ref{prop:tilts}.  Let \(N_{\rm in}\) count
straddles wholly inside \(r_{i-k}\), and let \(N_{\rm out}\)
count straddles whose entire raw traces lie outside
\(r_{i+k}\) and inside the large cutoff.  In particular,
\emph{outside} refers to the trace, not to the region enclosed
by the arc.  These counts have disjoint supports.  Put
\[
 M=N_n-N_{\rm in}-N_{\rm out}\ge0,\qquad
 \dd Q_i=
   \frac{b^{N_{\rm in}+N_{\rm out}}}
        {\EE b^{N_{\rm in}+N_{\rm out}}}\,\dd\PP .
\]
Every omitted straddle counted by \(M\) supplies a passage
in one of a fixed number of shell bands between
\(r_{i-k}\) and \(r_{i+k}\), allowing fixed buffer shifts.
Indeed it is neither confined below the lower cutoff
nor confined above the upper one; its radial travel or
its crossing between opposite rays has substantial
diameter near one of those levels.  Distinct straddles
give disjoint subarcs.  The small-box construction used
for \eqref{angular:nest-increments} therefore bounds
\(M\) by a fixed sum of shell passage counts.

These counts have all fixed exponential moments under
\(Q_i\).  For a count near the inner edge, condition on
states outside an intermediate radius below the support
of the outer tilt.  The outer factor is then constant,
and the conditional law is the inner tilt with a
compatible exterior partition.  Its edge moment is
bounded by Proposition~\ref{prop:tilts}.  For a count
near the outer edge, reverse this conditioning and use
the inward version of that proposition.  Middle buffers
miss both supports and use
Proposition~\ref{prop:passage-moments} directly.  H\"older's
inequality over the fixed number of counts gives
\begin{equation}\label{angular:omitted-moments}
 \EE_{Q_i}e^{tM}\le C_t\qquad(t\ge0),
\end{equation}
uniformly in \(i,n\) and the exterior states.  This
argument uses conditional estimates, not independence
of the two tilts.

Between the two tilt supports, choose two separated
regular subbands.  Proposition~\ref{prop:screens}
constructs proper actual wall-to-wall paths of opposite
colors in these bands.  Their conditional probabilities
are bounded below uniformly after fixing states away
from their supports; the two tilt factors are measurable
there.  Applying the construction successively gives an
event \(E_i\) with
\[
 Q_i(E_i)\ge p>0 .
\]
On \(E_i\), at least one raw straddle lies between the
two paths.  Indeed raw strands cannot cross either
actual path.  Along each flank, the end vertices of
the two paths have opposite colors, so the wall segment
between them contains an odd number of raw wall ports.
In the intervening disk, pairings of ports on a single
flank remove them in pairs.  Hence at least one raw
arc joins the two flanks.  It is trapped between the
blockers and is a localized straddle.

The full tilt \(\nu\) by \(b^{N_n}\) is obtained from
\(Q_i\) by the density
\(b^M/\EE_{Q_i}b^M\).  Cauchy--Schwarz yields
\[
 Q_i(E_i)^2
 \le \EE_{Q_i}[b^M\one_{E_i}]\,
       \EE_{Q_i}[b^{-M}\one_{E_i}],
\]
and consequently
\begin{equation}\label{angular:tilt-event-transfer}
 \nu(E_i)\ge
 \frac{p^2}
      {\EE_{Q_i}b^M\,\EE_{Q_i}b^{-M}}
 \ge c_b>0 .
\end{equation}
The last inequality follows from
\eqref{angular:omitted-moments}; because \(M\ge0\),
one of the two denominator factors is at most one.
Thus the same argument works for \(b<1\) and \(b>1\).

Choose \(c_1\log n-O(1)\) scales whose localization
bands are disjoint.  Their forced straddles are distinct,
so, pointwise,
\(N_n\ge\sum_i\one_{E_i}\).  Taking expectation under
the single full tilt and using
\eqref{angular:tilt-event-transfer} gives
\[
 \EE_{\nu}N_n\ge c_2\log n-O(1).
\]
For the finite graph,
\(\psi(t)=\log\EE e^{tN_n}\) is differentiable and
convex, with
\(\psi'(\log b_1)=\EE_{\nu}N_n\).
The supporting-line inequality between
\(\log b_1\) and \(\log b_2\) proves
\eqref{angular:strict-weight-gap}.

Finally choose any fixed \(0<b_1<d\) and put \(b_2=d\).
Since the summand for \(N_n=0\) has weight one,
\[
 \EE b_1^{N_n}\ge\PP(N_n=0)
       \ge\PP(\mathcal A_{Cn}).
\]
Lemma~\ref{angular:arm-straddle} compares the final
probability with \(\PP(\mathcal A_n)\) by a fixed factor.
Combining this with \eqref{angular:strict-weight-gap}
proves \eqref{angular:strict-flat-gap}.
\end{proof}

The costs $P,D_*$ enter the charge comparison of
Theorem~\ref{thm:charge-budget}.  The finite matrix argument next isolates
the two decay powers.  Section~\ref{sec:physical} proves their staggered
physical instance and then uses it in the scalar estimate of
Theorem~\ref{thm:physical-transfer} and the geometric application in
Corollary~\ref{cor:finite-extraction-rates}. The remaining arguments in
this section concern finite matrices alone.

\subsection{Finite-circle rows and the scalar normalization}

For the rest of this Section, temporarily allow the larger parameter range
\begin{equation}\label{angular:parameters}
 0<\lambda<\frac\pi2,\qquad d=2\cos\lambda>0,\qquad
 \rho=\frac\lambda\pi,\qquad v=\frac\pi2.
\end{equation}
The target range is $\pi/3<\lambda<\pi/2$.  No infinite-volume state is
used in this part of the argument.  Every row acts on a finite tensor
product, and every limit below first sends its circumference to infinity
with the entire finite row arrangement fixed.

In the ordered spin basis $++,+-,-+,--$, set
\begin{equation}\label{angular:R}
 R(u)=
 \begin{pmatrix}
 a(u)&0&0&0\\
 0&b(u)&1&0\\
 0&1&b(u)&0\\
 0&0&0&a(u)
 \end{pmatrix},\qquad
 a(u)=\frac{\sin(\lambda-u)}{\sin\lambda},\qquad
 b(u)=\frac{\sin u}{\sin\lambda}.
\end{equation}
Also write
\begin{equation}\label{angular:acr}
 c(u)=\frac{\sin(\lambda+u)}{\sin\lambda},\qquad
 r(u)=a(u)c(u)=1-b(u)^2.
\end{equation}
The scalar used to normalize a row is defined on the outer spectral
strip, independently of any Perron eigenvector.  The use of an inversion
relation together with symmetry and analyticity follows the classical
method developed by Baxter~\cite{Baxter1982Inversion}.  We define the
scalar by its boundary data and prove the identities needed here.

\begin{lemma}[Outer-strip scalar]\label{angular:scalar-normalizer}
There is a unique nonvanishing holomorphic function $F$ on
$0<\Re u<\lambda$ whose logarithmic modulus is the strip Poisson
integral of
\[
 h(y)=\tfrac12\log r(\ii y)
     =\tfrac12\log\left(1+\frac{\sinh^2 y}{\sin^2\lambda}\right)
\]
on both boundary lines and for which $F(\lambda/2)>0$.
It extends locally across both boundary lines, satisfies
\begin{equation}\label{angular:F-symmetries}
 F(\bar u)=\overline{F(u)},\qquad
 F(\lambda-u)=F(u),\qquad
 F(u)F(-u)=r(u),
\end{equation}
and has $F(0)=F(\lambda)=1$.  The last identity continues $F$
nonvanishingly to $-\lambda<\Re u<0$, and symmetry continues it to
$\lambda<\Re u<2\lambda$.  On the common domain,
\begin{equation}\label{angular:F-fusion}
 F(u)F(\lambda+u)=r(u).
\end{equation}
In the closed original strip, and in each closed substrip of
$-\lambda<\Re u<2\lambda$,
$\log|F(u)|-|\Im u|$ is bounded, with a bound depending on
$\lambda$ and the chosen substrip.
\end{lemma}

\begin{proof}
For $0<x<\lambda$, the two strip Poisson kernels are
\[
 P_0(x,t)=\frac{\sin(\pi x/\lambda)}
 {2\lambda\{\cosh(\pi t/\lambda)-\cos(\pi x/\lambda)\}},\qquad
 P_\lambda(x,t)=\frac{\sin(\pi x/\lambda)}
 {2\lambda\{\cosh(\pi t/\lambda)+\cos(\pi x/\lambda)\}}.
\]
Their masses are $1-x/\lambda$ and $x/\lambda$, respectively.  Their
sum has a uniformly bounded first absolute moment, and exponentially
decaying tails.  Thus
\[
 U(x+\ii y)=\int_\R
       \{P_0(x,y-t)+P_\lambda(x,y-t)\}h(t)\dd t
\]
is harmonic and $U(x+\ii y)-|y|$ is bounded, since $h(t)-|t|$ is
bounded.  A harmonic conjugate on the simply connected strip defines
$F=\exp(U+\ii V)$.  Fix its constant phase by $F(\lambda/2)>0$.
This also proves uniqueness.

The boundary data are real analytic.  For example, near any point of
the imaginary edge choose the analytic branch of
$\tfrac12\log r(u)$ that is real on that edge.  The difference of its
real part and $U$ vanishes there, so odd harmonic reflection extends
that difference across the edge.  This gives a local holomorphic
extension of $F$, still nonvanishing.  The other edge is identical.
Symmetries of the Poisson integral and the phase at $\lambda/2$ give
the first two identities in \eqref{angular:F-symmetries}.
On the imaginary edge,
$F(\ii y)F(-\ii y)=|F(\ii y)|^2=r(\ii y)$, so analytic
continuation gives the third identity.

The only zeros of $r$ closest to the imaginary axis are at
$u=\pm\lambda$; all its zeros are real.  Hence $r(u)/F(-u)$
extends $F$ without zeros to $-\lambda<\Re u<0$.  Reflection about
$\lambda/2$ gives the extension across the other edge and
\eqref{angular:F-fusion}.  Along the real interval $[0,\lambda]$,
$F$ is real, nonzero, and positive at the midpoint, hence positive
throughout.  At zero the third identity gives $F(0)^2=r(0)=1$;
the endpoint values follow.
The growth assertion on the extended substrips follows from
$F(-u)=r(u)/F(u)$ and symmetry, because
$\log|r(u)|-2|\Im u|$ is bounded at the two imaginary infinities
on every bounded real interval.  On a compact remaining rectangle
the extended function is continuous and nonzero.
\end{proof}

Let $\mathcal H_k=(\C^2)^{\otimes k}$, with its usual Hilbert norm,
and let $p$ denote an auxiliary spin.  For a $2$-by-$2$ matrix $D$,
define
\begin{equation}\label{angular:rows}
 T_k^D(u)=\Tr_p\bigl(D_pR_{p1}(u)\cdots R_{pk}(u)\bigr),\qquad
 A_k^D(z)=\frac{T_k^D(\rho(v-z))}{F(\rho(v-z))^k}.
\end{equation}
The auxiliary trace is unnormalized.  Subscripts $k$ will sometimes be
suppressed.  For $|\zeta|=1$ put
\begin{equation}\label{angular:twist-rows}
 K(\zeta)=\begin{pmatrix}\zeta&0\\0&\zeta^{-1}\end{pmatrix},
 \qquad P_{\zeta,k}=A_k^{K(\zeta)}.
\end{equation}
Choose $\epsilon\in\{-1,1\}$ and unitary twists satisfying
\begin{equation}\label{angular:twist-ratio}
 \frac\zeta\xi=-e^{\ii\epsilon\lambda}.
\end{equation}
There are two permitted marker conventions:
\begin{equation}\label{angular:marker-convention}
 (D,\tau)=\bigl(|-\rangle\langle+|,\epsilon\bigr)
 \quad\hbox{or}\quad
 (D,\tau)=\bigl(|+\rangle\langle-|,-\epsilon\bigr).
\end{equation}
For either convention write
\begin{equation}\label{angular:MN}
 M_k(z)=e^{\ii\tau\rho z}A_k^D(z),\qquad
 N_k(z)=M_k(z)P_{\zeta,k}(z)+P_{\xi,k}(z)M_k(z).
\end{equation}

\begin{theorem}[Homogeneous angular sandwich]\label{thm:angular-transfer}
For every fixed $\lambda$ in \eqref{angular:parameters}, let
$B_{\mathrm{out},k}$ and $B_{\mathrm{in},k}$ be finite products of
$P_{\xi,k}$ and $P_{\zeta,k}$, respectively, at real arguments in
$[-v,v]$.  Suppose that each product contains at least $n\ge1$ rows
at argument zero.  Then
\begin{align}
 \limsup_{k\to\infty}\sup_{x\in[-v,v]}
 \|B_{\mathrm{out},k}N_k(x)B_{\mathrm{in},k}\|
 &\le C n^{-1+\rho},\label{angular:sandwich-first}\\
 \limsup_{k\to\infty}\sup_{x,y\in[-v,v]}
 \|B_{\mathrm{out},k}
   \{N_k(x)-e^{\ii\tau(x-y)}N_k(y)\}B_{\mathrm{in},k}\|
 &\le C n^{-1-\rho}.\label{angular:sandwich-second}
\end{align}
Immediately to either side of each $N_k$, one may insert an even
number, bounded by a fixed $L$, of additional rows with the
corresponding twist, at arbitrary real arguments in $[-v,v]$.
The two terms in \eqref{angular:sandwich-second} may use different
additional lists.  The constant can then depend on $L$, as well as
on $\lambda$, but not on $n$ or on the real arguments.
Moreover,
\begin{equation}\label{angular:contraction}
 \limsup_{k\to\infty}\sup_{x\in[-v,v]}
       \|P_{\zeta,k}(x)\|\le1.
\end{equation}
The limit superior is taken for each fixed finite setup: its row
counts and any geometric parameters are fixed before $k\to\infty$.
\end{theorem}

We call \eqref{angular:sandwich-first} the \emph{low-stencil}
(amplitude) bound and \eqref{angular:sandwich-second} the
\emph{high-stencil} (phase-cancelled) bound.  The latter removes the
leading angular factor $e^{\ii\tau x}$; these names concern angular
cancellation, not the spectral cutoff used below.

We establish the theorem in finite-dimensional steps.  The analytic
argument only uses scalar normal families obtained from these finite
matrices; no common infinite-volume spectral representation is needed.

\begin{lemma}[Finite row identities]\label{angular:basic-rows}
For every $k$ and every unitary twist,
\begin{equation}\label{angular:normal-rows}
 [P_{\zeta,k}(z),P_{\zeta,k}(w)]=0,\qquad
 P_{\zeta,k}(z)^*=P_{\zeta,k}(-\bar z).
\end{equation}
In particular, this holomorphic family has a common orthonormal
eigenbasis.  For $|\Re z|\le v$,
\begin{equation}\label{angular:row-bound}
 \|A_k^D(z)\|\le2\|D\|.
\end{equation}
\end{lemma}

\begin{proof}
Let $J=\left(\begin{smallmatrix}0&1\\1&0\end{smallmatrix}\right)$.
Direct multiplication gives
\begin{align}
 R_{12}(x-y)R_{13}(x)R_{23}(y)
 &=R_{23}(y)R_{13}(x)R_{12}(x-y),\label{angular:YB}\\
 R(u)R(-u)&=r(u)I,\label{angular:R-inversion}\\
 J_pR(u)^{t_p}J_p&=R(\lambda-u),\label{angular:R-crossing}
\end{align}
where $t_p$ is transpose only in the auxiliary spin.  For completeness,
the nonconstant entries in the one-minus-spin sector of
\eqref{angular:YB} reduce to
\begin{align*}
 b(x-y)+a(x)b(y)&=b(x)a(y),\\
 b(y)+a(x)b(x-y)&=b(x)a(x-y),\\
 a(x)+b(x-y)b(y)&=a(x-y)a(y).
\end{align*}
These are sine addition identities.  The all-equal-spin sectors are
scalar, and spin reversal gives the remaining sector.  Equations
\eqref{angular:R-inversion} and \eqref{angular:R-crossing} follow
from their displayed $2$-by-$2$ blocks.

Slide the intertwiner $R_{12}(x-y)$ across the two auxiliary rows
using \eqref{angular:YB}.  Their boundary tensor
$K(\zeta)\otimes K(\zeta)$ commutes with the intertwiner because
$R$ preserves total spin.  For generic arguments the intertwiner is
invertible, and cyclicity of the auxiliary trace proves row
commutativity.  Analytic continuation removes the genericity
restriction.

On taking an adjoint, the order of the site factors reverses.
Auxiliary partial transpose reverses that order once more: its matrix
entries belonging to distinct physical sites commute.  Crossing then
gives, for any $D$,
\[
 T_k^D(u)^*=T_k^{J\bar D J}(\lambda-\bar u).
\]
For a unitary diagonal twist $J\overline{K(\zeta)}J=K(\zeta)$.
The symmetries of $F$ prove the second identity in
\eqref{angular:normal-rows}.  Commutativity with the adjoints makes
each row normal and gives simultaneous orthogonal diagonalization.

On $u=\ii y$, the real-coefficient formula for $R$ and inversion
show that $R(\ii y)/\sqrt{r(\ii y)}$ is unitary.  A product of
these matrices on the auxiliary and physical spins is unitary too.
Since $|F(\ii y)|=\sqrt{r(\ii y)}$ and auxiliary dimension is two,
its partial trace with $D$ has norm at most $2\|D\|$.  On
$u=\lambda+\ii y$, crossing converts the traced row into a reversed
product at $-\ii y$, with the transpose and spin flip moved to $D$.
The same bound follows.  For each fixed $k$, the growth of $F$ and
the explicit trigonometric entries of $R$ show that the normalized
matrix entries are bounded throughout the strip.  The bounded-strip
maximum principle, applied to every unit-vector matrix element,
gives \eqref{angular:row-bound}.
\end{proof}

\subsection{The fusion estimate}

The essential gain is a bound on the three-dimensional auxiliary
quotient in a two-row product.  It remains valid when $d<1$.
Fusion of shifted fundamental $R$-matrices at a degenerate auxiliary
channel is an established construction; Kulish, Reshetikhin, and
Sklyanin develop it in the rational $\mathrm{GL}(N)$ setting
\cite[Section~2, equations~(12)--(18)]{KulishReshetikhinSklyanin1981}.
The trigonometric quotient and its norm estimate below are proved directly.

\begin{lemma}[A spin-one quotient]\label{angular:fusion}
Use two auxiliary spins, labeled $1,2$, and one physical spin $j$,
and set
\begin{equation}\label{angular:G}
 G_j(u)=R_{1j}(u)R_{2j}(\lambda+u).
\end{equation}
The auxiliary line spanned by
$s_0=|+-\rangle+|-+\rangle$ is invariant, and $G_j$ acts there
as $r(u)I$ on the physical spin.  In the three-dimensional quotient,
use the ordered basis
\[
 e_+=|++\rangle,\qquad e_0=\sqrt d\,[|+-\rangle],\qquad
 e_-=|--\rangle,
 \qquad [|-+\rangle]=-[|+-\rangle].
\]
The quotient of $G_j/r$ is
$Y_j(u)=(b(u)/r(u))\mathcal L_j(u)$, where, in the ordered basis
$e_+\otimes+,e_+\otimes-,e_0\otimes+,e_0\otimes-,
e_-\otimes+,e_-\otimes-$,
\begin{equation}\label{angular:fused-blocks}
 \mathcal L(u)=
 \begin{pmatrix}
 -a&0&0&0&0&0\\
 0&c&\sqrt d&0&0&0\\
 0&-\sqrt d&-b&0&0&0\\
 0&0&0&-b&\sqrt d&0\\
 0&0&0&-\sqrt d&c&0\\
 0&0&0&0&0&-a
 \end{pmatrix}.
\end{equation}
In this fixed quotient Hilbert norm, for $|\Re u|<\lambda/2$,
\begin{equation}\label{angular:fusion-bound}
 \left\|\left.\frac{\prod_{j=1}^kG_j(u)}{r(u)^k}
                   \right|_{\mathrm{quotient}}\right\|
 \le C\left|\tan\frac{\pi u}{2\lambda}\right|^k.
\end{equation}
Constants allow the fixed change of auxiliary basis, but are
independent of $k$.
\end{lemma}

\begin{proof}
In the one-minus-spin sector, in the order
$(++-),(+-+),(-++)$, direct multiplication gives
\begin{equation}\label{angular:G-three}
 G(u)=
 \begin{pmatrix}
 bc&b&-b\\
 a&ac&0\\
 c&1&-b^2
 \end{pmatrix}.
\end{equation}
Because $ac=1-b^2=r$, the vector $(0,1,1)$ is an eigenvector
with eigenvalue $r$.  Spin reversal gives the other physical-spin
fiber of the invariant auxiliary line.  On the quotient, using
$a-c=-db$, the mixed block in the unrescaled basis is
\[
 b\begin{pmatrix}c&1\\-d&-b\end{pmatrix}.
\]
Rescaling its middle auxiliary vector by $\sqrt d$ produces the
first mixed block of \eqref{angular:fused-blocks}; spin reversal
negates that middle vector and produces the other block.  The two
unmixed states have value $-ab$.  This proves the complete quotient
formula.

On $\Re u=-\lambda/2$ one has
\[
 \bar c=-b,\qquad |a|^2=|b|^2+d,\qquad |b|=|c|.
\]
Each mixed block of $\mathcal L$ has adjoint-square $|a|^2I$;
the unmixed entries have the same modulus.  Since $r=ac$,
$Y=(b/r)\mathcal L$ is unitary on that boundary line.

On $\Re u=\lambda/2$, take the partial transpose of
$\mathcal L$ in the rescaled auxiliary basis and, if desired,
exchange $e_+,e_-$.  Its two mixed blocks have diagonal entries
$-a,-b$ and opposite real off-diagonal entries $\sqrt d,-\sqrt d$;
the two unmixed entries equal $c$.  Here
\[
 a=\bar b,\qquad |c|^2=|b|^2+d,
\]
so the transformed matrix is unitary times $|c|$.  Again the
factor $b/r$ makes its norm one.

Let $Q_k(u)=Y_1(u)\cdots Y_k(u)$.  On the left boundary
$\|Q_k\|=1$.  On the right boundary its auxiliary partial transpose
is the reversed product of the individually transposed factors,
because entries at distinct physical sites commute.  That product
is unitary.  Undoing a partial transpose on a three-dimensional
auxiliary space costs at most the fixed factor three in operator
norm: writing unit vectors as $\sum_{i=1}^3e_i\otimes x_i$ and
$\sum_{j=1}^3e_j\otimes y_j$, bound each transposed block by the
original operator norm and use
$\sum_i\|x_i\|\sum_j\|y_j\|\le3$.
Thus $\|Q_k\|\le3$ on the right boundary.

There are no zeros of $r$ in the closed fusion strip
$|\Re u|\le\lambda/2$.  Each quotient factor contains $b(u)$,
so $Q_k$ has a zero of order at least $k$ at zero.  The map
\[
 B(u)=\tan\frac{\pi u}{2\lambda}
\]
has its only zero in the strip at zero, has modulus one on its
boundary lines, and maps its interior into the unit disk.  Hence
$Q_k/B^k$ has a removable singularity at zero and no other
singularities there.  Its entries are bounded at both imaginary
infinities: $b\mathcal L/r$ is bounded there and $B$ tends to a
number of modulus one.  The strip maximum principle applied to
its scalar matrix elements gives \eqref{angular:fusion-bound}.
\end{proof}

\begin{lemma}[Normalized two-row inversion]\label{angular:two-row}
For $u=\rho(v-z)$ with $|\Re u|<\lambda/2$, put
$\mathfrak b(u)=|\tan(\pi u/(2\lambda))|$.  Then
\begin{align}
 \|P_{\zeta,k}(z)P_{\zeta,k}(z-\pi)-I\|
 &\le C\mathfrak b(u)^k,\label{angular:inversion}\\
 \|M_k(z)P_{\zeta,k}(z-\pi)
       +P_{\xi,k}(z)M_k(z-\pi)\|
 &\le C|e^{\ii\tau\rho z}|\mathfrak b(u)^k.
       \label{angular:marker-inversion}
\end{align}
The continued rows are those defined by Lemma~\ref{angular:scalar-normalizer}.
The errors vanish uniformly on every compact subset of this fusion
strip.  Equation~\eqref{angular:contraction} holds; more generally,
the limiting norm is at most one uniformly on compact subsets of
the closed physical strip $|\Re z|\le v$.
\end{lemma}

\begin{proof}
The product of the two auxiliary row traces can be regrouped as
the two-auxiliary trace of $\prod_jG_j(u)$.  Indeed the factors
$R_{1i}$ and $R_{2j}$ commute for $i\ne j$, so no same-site order
is changed.  The scalar denominator is $r(u)^k$ by
\eqref{angular:F-fusion}.

The twist $K(\zeta)\otimes K(\zeta)$ preserves the invariant
line and acts on it as the identity.  In a basis beginning with
that line, both this twist and every $G_j$ are block upper
triangular.  The trace therefore equals the invariant-line
contribution $I$ plus the trace of the quotient product against
a bounded auxiliary matrix.  The off-diagonal block contributes
exactly zero to the trace, regardless of its size.
Lemma~\ref{angular:fusion} proves \eqref{angular:inversion}.

For the marker use
\[
 Q=D\otimes K(\zeta)
       +e^{-\ii\tau\lambda}K(\xi)\otimes D.
\]
If $D=|-\rangle\langle+|$ and $\tau=\epsilon$, then
\[
 Qs_0=(\zeta^{-1}+e^{-\ii\epsilon\lambda}\xi^{-1})|--\rangle=0.
\]
If $D=|+\rangle\langle-|$ and $\tau=-\epsilon$, then
\[
 Qs_0=(\zeta+e^{\ii\epsilon\lambda}\xi)|++\rangle=0.
\]
Both cancellations use \eqref{angular:twist-ratio}.  Thus $Q$
kills the invariant line.  Its product with the block triangular
fusion product again has zero contribution from the off-diagonal
block, and now its invariant-line contribution is zero as well.
Taking its quotient trace and extracting $e^{\ii\tau\rho z}$
proves \eqref{angular:marker-inversion}.

On the right edge $z=v+\ii y$, one has $z-\pi=-\bar z$.
Consequently \eqref{angular:normal-rows} turns
\eqref{angular:inversion} into approximate unitarity.  On the left
edge apply the same identity at $z+\pi$.  The errors converge to
zero uniformly for $y$ in a fixed compact interval.  To pass into
the strip, truncate at $|\Im z|=Y$.  On the two vertical sides
the norm is at most $1+o_k(1)$; on the horizontal sides the
uniform bound is two.  The harmonic measure of the horizontal
sides tends uniformly to zero on any fixed compact subset as
$Y\to\infty$.  The maximum principle for unit-vector matrix
elements, first with $k\to\infty$ and then $Y\to\infty$, proves
the claimed limiting contraction.  The edge estimates include
the real endpoints.
\end{proof}

\subsection{Scalar limits and continuation through the strip edges}

Throughout this subsection, $0<\rho<1/2$ is fixed.  Every limit in the
number $k$ of sites is taken with the finite list of rows, its real
arguments, and all spectral cutoffs fixed.  In particular, estimates in
the damping parameter $n$ are obtained only after this limit.  Constants
may depend on $\rho$ and on the permitted bound for the number of
additional rows, but not on $k$ or $n$.

Put $\mathcal S=\{z:|\Re z|<v\}$, where $v=\pi/2$, and use
$w=e^{\ii z}$ to identify $\mathcal S$ with the right half-plane.  The
branch of $z=-\ii\log w$ in this identification is always the one with
$|\Re z|<v$.  For a fixed twist, all the rows have a common orthonormal
eigenbasis.  Their eigenvalues are therefore holomorphic functions of
$z$; the same basis diagonalizes their analytic continuations, by the
identity theorem.  The adjoint identity~\eqref{angular:normal-rows}
for the rows becomes
\begin{equation}\label{angular:eigenvalue-reflection}
 f_k(\overline w)=\overline{f_k(w)}
 \qquad (\Re w>0).
\end{equation}
In particular, $f_k(1)$ is real.  Here and below $P(0)$ refers to the
argument $z=0$, which corresponds to $w=1$.

\begin{lemma}[Scalar edge continuation]\label{angular:edge-continuation}
Fix $s>0$.  From any sequence of row eigenvalues satisfying
$|f_k(1)|\ge e^{-s}$ one can extract a subsequence converging locally
uniformly in the right half-plane to a holomorphic function $f$ that
is not identically zero.  It satisfies
\[
 |f|\le1,\qquad f(1)\in\R,\qquad |f(1)|\ge e^{-s}.
\]
Its only possible zeros in the right half-plane lie on the positive
real axis.  At every nonzero point of the imaginary axis, $f$ extends
holomorphically and without zeros through the axis, and has modulus
one there.  Along the chosen subsequence, the eigenvalues and their
inverses are locally bounded through each such boundary point.
\end{lemma}

\begin{proof}
The row bound~\eqref{angular:row-bound} gives a normal family in
$\mathcal S$.  The lower bound at $1$ prevents its limit from being
identically zero.  The asymptotic contraction
\eqref{angular:contraction} gives $|f|\le1$ throughout the strip: its
boundary version, together with the uniform row bound, can first be
applied on a finite rectangle, and the harmonic measure of the two
horizontal sides tends to zero as the rectangle grows.  Equivalently,
one can use the interior contraction already obtained from this
argument.  Equation~\eqref{angular:eigenvalue-reflection} passes to the
limit.

The inversion estimate~\eqref{angular:inversion} makes $f_k$ zero-free
on every compact subset of each spectral band about a strip edge, for
all sufficiently large $k$.  Under $u=\rho(v-z)$, the band
$|\Re u|<\lambda/2$ about the right edge is $0<\Re z<\pi$;
the corresponding band about the left edge is $-\pi<\Re z<0$.
These two bands cover the strict strip except for its center line.
Hurwitz's theorem therefore excludes
zeros of $f$ off that center line, which maps to the positive real
axis in the $w$-plane.

We give the additional argument needed at an edge; mere convergence
on the strict strip would not suffice.  Fix a disk $D$ centered at
$z_0=v+\ii y_0$, with closure compactly contained in the right-edge
spectral band.  In this paragraph write $f_k(z)$ for $f_k(e^{\ii z})$.
On $D$,
\[
 q_k(z):=f_k(z)f_k(z-\pi)\longrightarrow1
\]
uniformly.  For all sufficiently large $k$, choose the analytic square
root of $q_k$ near $1$, and set
\[
 g_k(z)=\frac{f_k(z)}{\sqrt{q_k(z)}},
 \qquad U_k(z)=-\log|g_k(z)|.
\]
Both eigenvalue factors are nonzero on $D$.  On its vertical diameter
the adjoint identity gives $q_k(z)=|f_k(z)|^2$, so that $|g_k(z)|=1$.
Consequently $U_k$ is harmonic on $D$ and odd under reflection in that
diameter.  On the physical half of $D$, the row bound and $q_k\to1$
give $U_k\ge-C$.  At a fixed point strictly inside that half-disk,
$U_k$ is bounded above as well: the limiting eigenvalue has no zero
there by the preceding paragraph.

For clarity, the elementary boundary estimate used here is as follows.
For a harmonic function $U$ odd across a diameter of a disk, its
Poisson formula is an integral over one semicircle with kernel
\[
 K(z,t)=P(z,t)-P(z,t^{\rm ref}),
\]
where $P$ is the disk Poisson kernel and $t^{\rm ref}$ is the reflected
boundary point.  This kernel is nonnegative when $z$ and $t$ are on
the same side of the diameter.  On any strictly smaller half-disk,
\begin{equation}\label{angular:half-disk-kernel}
 0\le K(z,t)\le C K(z_*,t)
\end{equation}
for a fixed interior radial point $z_*$ on that side.  The assertion
follows directly from the Poisson formula: the difference has the
product of the two signed distances to the diameter in its numerator,
and its remaining denominators are bounded above and below on the
smaller disk.  If $U\ge-C$ on the physical half-disk, the negative
part of its semicircle data has a bounded contribution.  The formula
at $z_*$ therefore bounds the positive contribution there, and
\eqref{angular:half-disk-kernel} bounds it at all points of the smaller
half-disk.  Reflection supplies the other half.

Applying this observation on an intermediate disk inside $D$ bounds
$|U_k|$ on a smaller closed disk.  Thus $f_k$ and $f_k^{-1}$ are
uniformly bounded there.  Normal families extend the already chosen
interior limit across the edge, and the limit has modulus one on the
diameter.  This extension is unique, so taking further subsequences
does not change it.  The left edge is treated with the shifted
inversion identity.  These two arguments cover every nonzero point
of the imaginary axis.
\end{proof}

\subsection{Almost constant low spectral modes}

The outer row products in Theorem~\ref{thm:angular-transfer} damp
eigenvalues whose modulus at zero stays below one.  We therefore first
analyze the complementary spectral subspace, on which that modulus is
at least $e^{-s}$.  The next lemma shows that its row eigenvalues remain
close to one of the two real signs over the full physical angle interval.

\begin{lemma}[Half-annulus and sign estimates]
\label{angular:half-annulus}
There are constants $c_0,s_0,C>0$ such that the following holds for
$0<s\le s_0$.  Define
\[
 \mathcal A_s=\{w:c_0^{-1}s<|w|<c_0/s\},
 \qquad
 \mathcal A_s^+=\overline{\mathcal A_s}\cap\{\Re w\ge0\}.
\]
Every scalar limit in Lemma~\ref{angular:edge-continuation} obeys
\begin{equation}\label{angular:scalar-annulus-bound}
 C^{-1}\le |f(w)|\le C\quad(w\in\mathcal A_s^+),
 \qquad
 f(w)=\sigma+O(s)\quad(|w|=1,\ \Re w\ge0),
\end{equation}
where $\sigma=\sgn f(1)$ and the boundary values use its continuation.
The same assertions hold, in operator norm and with a limsup as
$k\to\infty$, for the diagonal compressions
\[
 E_{\zeta,k}(s)=
 \one_{\{|P_\zeta(0)|\ge e^{-s}\}},
 \qquad
 p_{\zeta,k}(w)=P_\zeta(z)|_{\operatorname{ran}E_{\zeta,k}(s)}.
\]
For the sign assertion, restrict further to either sign of
$P_\zeta(0)$.  The conclusions also hold for any further spectral
restriction.  The compressed rows and their inverses are locally
bounded in neighborhoods of the two boundary rays of $\mathcal A_s$.
The same statements apply to the twist $\xi$.
\end{lemma}

\begin{proof}
Multiply $f$ by its sign at $1$ and put $a=|f(1)|$.  For a nonconstant
function, Schwarz--Pick in the right half-plane gives, for $r>0$,
\[
 \left|\frac{\sigma f(r)-a}{1-a\sigma f(r)}\right|
 \le t:=\frac{|r-1|}{r+1}.
\]
Writing the fraction on the left as $h$, one obtains
\[
 |1-\sigma f(r)|
 = (1-a)\left|\frac{1-h}{1+ah}\right|
 \le (1-a)\frac{1+t}{1-t}
 \le s\max(r,r^{-1}).
\]
Constant functions satisfy the same conclusion directly.  Choose
$c_0$ sufficiently small.  This estimate excludes zeros on the
positive real axis even on an annulus with inner radius
$c_0^{-1}s/4$ and outer radius $4c_0/s$.  Lemma
\ref{angular:edge-continuation} excludes zeros elsewhere in its right
half.  On that half-annulus, therefore,
\[
 U(w):=-\log|f(w)|
\]
is a nonnegative harmonic function, vanishing on the imaginary-axis
boundary segments and extending oddly through them.  On the positive
real axis the same estimate gives
\begin{equation}\label{angular:radial-log-bound}
 U(r)\le Cs\max(r,r^{-1}).
\end{equation}

Fix $r\in[c_0^{-1}s,c_0/s]$ and rescale by $r$.  A bounded Harnack
chain inside $\{1/4<|w|<4,\ \Re w>0\}$ controls $U$ on the unit
semicircle away from its endpoints by $CU(r)$.  Near either endpoint,
use a disk centered on the imaginary axis, of fixed radius after
rescaling.  Its interior reference point is controlled by the same
Harnack chain; the reflected Poisson estimate
\eqref{angular:half-disk-kernel}, now with nonnegative data on the
physical half, controls the smaller half-disk up to the diameter.
Consequently
\begin{equation}\label{angular:semicircle-log-bound}
 0\le U(w)\le Cs\max(r,r^{-1})
 \quad (|w|=r,\ \Re w\ge0).
\end{equation}
The right side is bounded by a fixed constant on the asserted
annulus, proving its uniform inverse bound.

For $r$ in a fixed neighborhood of $1$, the bound is $Cs$.
After odd reflection it bounds $|U|$ in fixed neighborhoods of the
two unit semicircle endpoints as well.  Harmonic gradient estimates
therefore give $|\nabla U|\le Cs$ in a neighborhood of the whole unit
semicircle.  Choose a branch of $\log(\sigma f)$ in a thin simply
connected neighborhood of that semicircle, with value $\log a$ at
$1$.  The Cauchy--Riemann equations identify the modulus of its
derivative with $|\nabla U|$.  Integration along the semicircle gives
$\log(\sigma f)=O(s)$, and hence the second assertion in
\eqref{angular:scalar-annulus-bound}.

We spell out why these scalar arguments are uniform for the growing
matrices.  If any compressed diagonal norm assertion failed, choose
an offending eigenvalue, a point in the indicated compact set, and a
subsequence of site numbers.  The points have a convergent
subsequence.  Lemma~\ref{angular:edge-continuation} supplies a scalar
limit, including local convergence through an edge when necessary;
\eqref{angular:scalar-annulus-bound} then contradicts the failure.
For the sign assertion one first fixes the sign subsequence.  The
same argument rules out zeros of the finite compressed eigenvalues
on the compact half-annulus for all sufficiently large $k$, so that
their inverses are defined there.  No bound on the number of
eigenvalues is used.

Finally, scalar continuation across an imaginary boundary interval
with modulus one is given by Schwarz reflection,
$f(w)=1/\overline{f(-\overline w)}$ on its other side.  The just-proved
inverse bounds therefore control this continuation near either ray,
as well as its inverse.  The edge-continuation argument and the same
subsequence contradiction give the corresponding local bounds for
the finite compressed matrices.  One may decrease $c_0$ once to keep
all neighborhoods inside the wider annulus used above.
\end{proof}

\subsection{Marker continuation and its two leading modes}

The diagonal row estimates do not directly bound a marker between the
two twist spaces.  The fusion identity continues its scalar matrix
elements around an annulus.  Their Laurent coefficients then separate
the leading angular mode from the remainder canceled by a high stencil.

\begin{lemma}[Annular marker continuation]
\label{angular:marker-continuation}
Let $E_\xi,E_\zeta$ be the low spectral projections from Lemma
\ref{angular:half-annulus}, possibly further spectrally restricted,
and put
\[
 m_k(w)=E_\xi M(z)E_\zeta,\qquad w=e^{\ii z}.
\]
For any sequence of unit vectors $a_k\in\operatorname{ran}E_\xi$ and
$b_k\in\operatorname{ran}E_\zeta$, one can extract a subsequence for
which $\langle a_k,m_k(w)b_k\rangle$ converges on the right half of
$\mathcal A_s$ to the restriction of a holomorphic function $m$ on
the full open annulus $\mathcal A_s$.  Convergence includes the closed
right half of every strictly smaller annulus, using continuation at
its imaginary boundary points.  The limit satisfies
\begin{equation}\label{angular:marker-growth}
 |m(w)|\le C|w|^{\tau\rho}\qquad(w\in\mathcal A_s).
\end{equation}
On the left half-annulus its values are the limits of the same
matrix elements of
\begin{equation}\label{angular:external-marker}
 m_k^{\rm ext}(-w)
 =-p_{\xi,k}(w)^{-1}m_k(w)p_{\zeta,k}(w)^{-1},
 \qquad \Re w>0.
\end{equation}
\end{lemma}

\begin{proof}
The row bound~\eqref{angular:row-bound}, including the scalar factor
in $M$, gives
\[
 \|m_k(w)\|\le C|w|^{\tau\rho}
\]
on the physical right half-plane.  Lemma
\ref{angular:half-annulus} gives the same limsup bound, up to a
constant, for \eqref{angular:external-marker}.  Thus the scalar
matrix elements of both expressions are normal families on their
respective half-annuli.  It remains to check that their limits glue
through both imaginary rays, including local boundedness there.

Use $z$ coordinates near $z=v+\ii y$.  Compressing the fusion
identity~\eqref{angular:marker-inversion} gives
\begin{equation}\label{angular:compressed-gluing}
 m_k(z)p_{\zeta,k}(z-\pi)
 +p_{\xi,k}(z)m_k(z-\pi)=o(1)
\end{equation}
uniformly on compact neighborhoods, in operator norm.  Here $m_k(z)$
abbreviates the continued compression at argument $z$, and the same
convention applies to $p$.  The error tends to zero for each fixed
$s$ and compact neighborhood.  The inversion identity also gives
\[
 p_{\xi,k}(z)=p_{\xi,k}(z-\pi)^{-1}+o(1),
\]
and the analogous relation for the other twist.  All these inverses
are locally bounded by Lemma~\ref{angular:half-annulus}.  In a small
neighborhood of the edge, one of $z,z-\pi$ is in the physical closed
strip.  Solving \eqref{angular:compressed-gluing} for the other marker
shows that both continued marker expressions are locally bounded.
Solving for $m_k(z)$ in particular gives
\[
 m_k(z)=
 -p_{\xi,k}(z-\pi)^{-1}m_k(z-\pi)
  p_{\zeta,k}(z-\pi)^{-1}+o(1).
\]
Since $e^{\ii(z-\pi)}=-e^{\ii z}$, this is exactly the compatibility
of the physical and external expressions across the upper imaginary
ray.  Across the lower ray use the same identity with $z+\pi$ in
place of $z$ and solve in the opposite direction.  Again one argument
is physical, and inversion produces precisely
\eqref{angular:external-marker}.

Montel's theorem and a diagonal subsequence over compact sets now
give the asserted holomorphic function on the full annulus.  This
uses scalar matrix elements only: the two twist spaces may have
different eigenbases and dimensions tending to infinity.  There is
no assertion that limits of individual matrix factors can be
multiplied.  Their finite-dimensional identities and uniform norm
bounds are used before taking subsequences.  Bound
\eqref{angular:marker-growth} follows on each half-annulus from its
defining expression, and on the rays by continuation.
\end{proof}

\begin{lemma}[Low-band estimates]\label{angular:low-band}
For $0<s\le s_0$, let $E_\xi,E_\zeta$ be as in Lemma
\ref{angular:marker-continuation}.  Uniformly for real $x,y\in[-v,v]$,
\begin{align}
 \limsup_{k\to\infty}\|E_\xi N(x)E_\zeta\|
 &\le Cs^{1-\rho},\label{angular:low-band-size}\\
 \limsup_{k\to\infty}
 \|E_\xi\{N(x)-e^{\ii\tau(x-y)}N(y)\}E_\zeta\|
 &\le Cs^{1+\rho}.\label{angular:low-band-difference}
\end{align}
The estimates persist if each occurrence of $N$ is multiplied,
immediately on each side separately, by a product of an even bounded
number of the corresponding twist rows at real arguments in $[-v,v]$.
Each of the four lists has even length; the lists for the two
occurrences may differ.
\end{lemma}

\begin{proof}
Expand the scalar limit in Lemma~\ref{angular:marker-continuation} as
$m(w)=\sum_{j\in\Z}b_jw^j$.  Cauchy's estimate on circles just inside
the two annular radii, with a fixed proportional margin, gives
\begin{equation}\label{angular:laurent-coefficients}
 |b_j|\le C(C's)^{|j-\tau\rho|}.
\end{equation}
For example, use the outer circle when $j>\tau\rho$ and the inner
circle when $j<\tau\rho$.  Shrinking $s_0$ makes $C's_0<1/2$.
Because $0<\rho<1/2$, summing these bounds on the unit circle yields
\begin{align}
 |m(w)|&\le Cs^\rho,\label{angular:marker-small}\\
 \frac{m(w)-m(-w)}2
 &=b_\tau w^\tau+O(s^{1+\rho}),
 &|b_\tau|&\le Cs^{1-\rho}.
 \label{angular:odd-mode}
\end{align}
Indeed $j=0$ has exponent $\rho$, the nearest nonzero exponent is
$1-\rho$, and among odd indices other than $\tau$ the smallest
exponent is $1+\rho$.  The remaining geometric tails are uniform.

By the arbitrary-unit-vector subsequence argument in Lemma
\ref{angular:marker-continuation}, \eqref{angular:marker-small} gives
\begin{equation}\label{angular:compressed-marker-small}
 \limsup_{k\to\infty}
 \sup_{|w|=1,\ \Re w\ge0}\|m_k(w)\|\le Cs^\rho.
\end{equation}
The supremum causes no difficulty: choose offending points, pass to
a convergent subsequence on the compact semicircle, and use the
local uniform convergence, including at its endpoints.

Further split each low spectral projection according to the sign of
$P(0)$, and first work on one pair of sign blocks, with signs
$\sigma_\xi,\sigma_\zeta\in\{1,-1\}$.  Lemma
\ref{angular:half-annulus} gives, in limsup uniformly on the physical
unit semicircle,
\[
 p_{\xi,k}=\sigma_\xi I+O(s),\qquad
 p_{\zeta,k}=\sigma_\zeta I+O(s),
\]
and the same estimates for their inverses.  Combining
\eqref{angular:external-marker} and
\eqref{angular:compressed-marker-small}, every scalar subsequential
limit satisfies
\begin{equation}\label{angular:marker-parity}
 m(-w)=-\sigma_\xi\sigma_\zeta m(w)+O(s^{1+\rho}).
\end{equation}
The compressed anticommutator is exactly
\[
 E_\xi N(z)E_\zeta
 =m_k(z)p_{\zeta,k}(z)+p_{\xi,k}(z)m_k(z).
\]
If the signs disagree, its leading terms cancel and it is
$O(s^{1+\rho})$ in limsup.  If both signs equal $\sigma$,
\eqref{angular:marker-parity} and \eqref{angular:odd-mode} instead give,
in each tested scalar subsequential limit,
\begin{equation}\label{angular:leading-mode}
 E_\xi N(x)E_\zeta
 =2\sigma b_\tau e^{\ii\tau x}+O(s^{1+\rho}).
\end{equation}
This equation denotes the corresponding scalar matrix element, not
an identification of spaces of different dimensions.  Its leading
term has size $O(s^{1-\rho})$ and cancels between $x$ and $y$ after
multiplication by $e^{\ii\tau(x-y)}$.  Applying the same argument to
arbitrary sequences of test vectors and of $x,y$ proves
\eqref{angular:low-band-size}--\eqref{angular:low-band-difference}.
Summing the four sign blocks changes the constant only.

Finally, an even product of at most a fixed number of twist rows
equals $I+O(s)$ on its sign block, because every factor equals
$\sigma I+O(s)$.  Multiplying the compressed anticommutator by such
products changes it by at most $Cs\|m_k\|$ in limsup, hence by
$Cs^{1+\rho}$.  This reasoning is separate for the two terms, so their
additional lists need not agree.  All these additional rows commute
with the spectral projections on their own side.  This proves the
last assertion.
\end{proof}

\subsection{Damping the remaining spectral bands}

\begin{proof}[Proof of Theorem~\ref{thm:angular-transfer}]
The row contraction is \eqref{angular:contraction}.  It remains to
insert the two outer products in the low-band estimates.  For either
twist write
\[
 h_k=-\log|P(0)|,
\]
with value $+\infty$ on the kernel.  The spectral projections of $h_k$
commute with every row of their twist.  If a band has positive lower
endpoint $t$, the at least $n$ zero-argument factors in its outer
product give the bound $e^{-nt}$ on that band.  The other factors
have asymptotic norm at most one.  Since their list is finite and
fixed while $k\to\infty$, the full product has limsup norm at most
$e^{-nt}$ there.  On a band without a positive lower endpoint its
limsup norm is at most one by contraction of all the factors.

For $n$ large enough, choose $J$ so that
$2^J/n\le s_0<2^{J+1}/n$, and split each side into
\[
 I_0=(-\infty,1/n],\qquad
 I_i=(2^{i-1}/n,2^i/n]\quad(1\le i\le J),
 \qquad I_\infty=(2^J/n,+\infty].
\]
Set $D_0=1$ and $D_i=e^{-2^{i-1}}$ for $i\ge1$.
For a pair of finite bands $I_i,I_j$, apply Lemma
\ref{angular:low-band} with
\[
 s=\frac{\max(2^i,2^j)}n.
\]
Both band projections are further restrictions of the appropriate
low projections.  Their outer products contribute $D_iD_j$.
For either exponent $a\in\{1-\rho,1+\rho\}$, the triangle inequality
therefore bounds the sum of these blocks by
\begin{equation}\label{angular:dyadic-sum}
 C n^{-a}\sum_{i,j=0}^{J}
 D_iD_j\max(2^i,2^j)^a
 \le C_a n^{-a}.
\end{equation}
The infinite double sum converges, since the $D_i$ decrease
exponentially in $2^i$.  In particular there is no factor counting
the number of bands.

The remainder has $h_k>2^J/n>s_0/2$ and therefore contributes
$e^{-ns_0/2}$ on its side.  Group all blocks with a remainder on the
left into a single compression, and then all remaining blocks with
a remainder on the right into another.  The unrestricted row bounds
control $N$ and its two-term difference; contraction controls the
outer product on the other side.  Thus these two groups together
contribute at most $Ce^{-ns_0/2}$ in limsup.  Additional bounded row
lists cause only a fixed constant, and their low-band contribution
was already covered by Lemma~\ref{angular:low-band}.

Taking $a=1-\rho$ in \eqref{angular:dyadic-sum} proves the first
asserted sandwich bound.  Taking $a=1+\rho$ proves the second, with
its phase $e^{\ii\tau(x-y)}$ and with independently chosen additional
even row lists.  The bounded range of smaller $n$ follows directly
from contraction of the outer products and the uniform marker bound,
after enlarging the constant.  At every stage
the band decomposition is finite before taking the site-number
limsup; only the resulting scalar estimates are then summed and
bounded uniformly in $n$.
\end{proof}
